% Fichier généré automatiquement par tools/generate_corrections.py.
% Source : PERF/PERF-02-Marges/62_Palettisation/62_Palettisation.tex
% Statut : complet (3/3 question(s) renseignée(s), 0 reprise(s)).
\begin{corrigebox}[Corrigé extrait de la banque]
\CorrigeQuestion{1}
On souhaite une marge de phase de 45\degres. On cherche donc $\omega_{\varphi}$ tel que 
$\varphi\left(\omega_{\varphi}\right)=-180+45 = -135\degres$.

$\varphi(\omega)=-90-\arg\left(1+\tau_m j \omega \right) =-90-\arctan\left(\tau_m  \omega \right)  $.

On a donc $\varphi\left(\omega_{\varphi}\right)=-135$
$\Leftrightarrow -90-\arctan\left(\tau_m  \omega_{\varphi}\right)  = -135 $
$\Leftrightarrow -\arctan\left(\tau_m  \omega_{\varphi} \right)  = -45 $
$\Leftrightarrow \arctan\left(\tau_m  \omega_{\varphi} \right)  = 45 $
$\Rightarrow \tau_m  \omega_{\varphi} = 1 $
$\Rightarrow \omega_{\varphi}  = \dfrac{1}{\tau_m}= \dfrac{1}{5\times 10^{-3}}$
$\Rightarrow \omega_{\varphi}  = \SI{200}{rad.s^{-1}}$.

Par suite, il faut que le gain soit nul en $\omega_{\varphi}$.

On a donc $\indice{G}{dB}(\omega)=20\log k_{BO}-20\log\omega-20\log\sqrt{1+\omega^2\tau_m^2}$.
En $\omega_{\varphi} =  \dfrac{1}{\tau_m}$ :
$\indice{G}{dB}(\omega_{\varphi})=0$ 
$\Leftrightarrow 20\log k_{BO}-20\log \dfrac{1}{\tau_m}-20\log\sqrt{1+\dfrac{1}{\tau_m^2}\tau_m^2}=0$
$\Leftrightarrow \log k_{BO}+\log \tau_m-\log\sqrt{2}=0$
$\Leftrightarrow \log \dfrac{k_{BO} \tau_m}{\sqrt{2}}=0$
$\Leftrightarrow \dfrac{k_{BO} \tau_m}{\sqrt{2}}=1$
$\Leftrightarrow k_{BO}=\dfrac{\sqrt{2}}{\tau_m}$.

(\textbf{A vérifier})
$k_{BO}=282,8$.
\CorrigeQuestion{2}
$k_{BO}=\dfrac{k_c k_m k_r}{N}$; donc 
$k_c=\dfrac{Nk_{BO}}{ k_m k_r} = \dfrac{200 \times 282,8 }{4\times 30}=471$.
\CorrigeQuestion{3}
Il y a une intégration dans la correcteur. La FTBO est de classe 1 est le système est précis en position.

\begin{enumerate}
  \item $k_{BO}=\dfrac{\sqrt{2}}{\tau_m}$.
  \item $k_c=\dfrac{\sqrt{2}N}{\tau_m k_m k_r} = 471,1$.
  \item $\varepsilon_s=0$.
\end{enumerate}
\end{corrigebox}
